\section{Discussion}
\label{sec:discussion}

\paragraph{Why we still think the VLM route is correct.}
Our main results are unflattering but the broader argument of the
paper stands: a parkour trick name is a \emph{linguistic} label over
a set of physical motions, not a set of motions in itself. ``Cork''
is a word that athletes and judges agreed means ``a backflip with a
half twist on a tilted axis''. It is a compositional entry in a
shared taxonomy that VLMs have absorbed from the parkour and tricking
internet at scale. A trained video classifier has to discover that
taxonomy from a few hundred clips per family and---as we show in
Sec.~\ref{sec:journey}---consistently fails to do so on our dataset.
The finding from our benchmark is not ``VLMs don't know parkour''
but rather ``frontier VLMs know parkour nouns better than they count
parkour revolutions, and in headless automation mode they have a
tendency to answer the wrong noun anyway.''

\paragraph{The automation gap.}
The most surprising and practically actionable result of the paper
is the gap between interactive and headless behaviour of the same
VLM on the same clips. In the absence of a clean explanation we
cannot, in good conscience, claim that a one-shot headless pipeline
is a solved problem. But we can name the gap and invite follow-up
work. Hypotheses worth testing in future work include:
\begin{itemize}
  \setlength\itemsep{-1pt}
  \item Interactive chat sessions carry a richer system prompt that
        includes instructions to reason carefully before answering.
  \item Headless mode's default temperature produces more
        confident-but-wrong samples than a human user would see,
        because the human would retry or ask a follow-up question.
  \item Tool-call accounting: interactive mode may allow the model
        more internal scratchpad reasoning before producing a final
        answer.
  \item Sampling variance: with 13 clips and one shot per model,
        every run is a draw from a noisy distribution.
\end{itemize}
Each of these is a small empirical study that would sharpen the
present result. We explicitly position this paper as a baseline that
the next wave of experiments can improve on.

\paragraph{The matcher is load-bearing, even when the predictions are
wrong.}
Table~\ref{tab:matcher_levels} shows that whenever a VLM returns a
trick name, the fuzzy matcher snaps it to a FIG entry with high
confidence, usually at the L1 exact level. The cascade therefore
carries very little of the load on our benchmark---the VLMs happen
to emit in-table names, even when those names are wrong. We still
believe the cascade is a load-bearing piece of the system because
(i)~it is the component that makes the VLM interface deterministic
and auditable, and (ii)~we saw informal VLM queries produce names
like ``back double full'', ``cork screw'', and ``krok'' that do need
the L2--L4 levels. A longer benchmark would likely exercise the
full cascade more thoroughly.

\paragraph{What the failure gallery tells us about the next system.}
All three of the recurring errors documented in
Sec.~\ref{sec:experiments:failures}---twist under-counting, takeoff
confusion, wall-context loss---are \emph{physics} errors that a
classical pose-tracking pipeline would be well-positioned to fix if
its input were reliable. That observation suggests a very concrete
next architecture: use the VLM only for the trick \emph{noun} (what
family of move is this), and delegate the \emph{attributes} (flip
count, twist count, direction, takeoff) to a deterministic physics
component driven by a 3D pose estimator fine-tuned on acrobatic
motion. The FIG table is compositional, so two small, well-aligned
signals would cover the whole space without the VLM having to count
twists at all. We leave this hybrid to future work; the present
paper argues only that the noun-only VLM path is the right starting
point and that the physics path remains necessary.

\paragraph{Cost, latency, deployability.}
Across the three configurations we benchmarked, a single trick
identification costs tens of seconds and a few cents or nothing
at all (Claude via an existing CLI subscription). A full competition
run of eight tricks is therefore fast and cheap enough to be useful
as an advisory tool during a live event, even if the accuracy
remains low: a human judge could in principle be shown the full
D-score card and correct only the entries the system flagged as
uncertain. None of our experiments tested this interactive
judge-assistance mode, but the system's text-level interfaces
(FIG name + match level + reasoning in every output) make it
straightforward to plug in.

\paragraph{Ethical note.}
We explicitly do not recommend using the raw \sys{} scorecard as a
substitute for a human judge's final score. The current accuracy is
too low, and the system is deliberately transparent about which
decisions it was confident in and which it was not. We hope the
fact that \sys{} is open-source and fully auditable---every
prediction records its model version, prompt hash, match level,
and raw text---means that any follow-up deployment will be on terms
a parkour federation can inspect before trusting.
